\section{Derived functors in general}
\label{section-derived-functors}

\noindent
A reference for this section is Deligne's expos\'e XVII in \cite{SGA4}.
A very general notion of right and left derived functors exists where
we have an exact functor between triangulated categories, a multiplicative
system in the source category and we want to find the ``correct'' extension
of the exact functor to the localized category.

\begin{situation}
\label{situation-derived-functor}
Here $F : \mathcal{D} \to \mathcal{D}'$ is an exact functor of triangulated
categories and $S$ is a saturated multiplicative
system in $\mathcal{D}$ compatible with the structure
of triangulated category on $\mathcal{D}$.
\end{situation}

\noindent
Let $X \in \Ob(\mathcal{D})$. Recall from
Categories, Remark \ref{categories-remark-left-localization-morphisms-colimit}
the filtered category $X/S$ of arrows $s : X \to X'$ in $S$ with source $X$.
Dually, in
Categories, Remark \ref{categories-remark-right-localization-morphisms-colimit}
we defined the cofiltered category $S/X$ of arrows $s : X' \to X$ in $S$
with target $X$.

\begin{definition}
\label{definition-right-derived-functor-defined}
Assumptions and notation as in
Situation \ref{situation-derived-functor}.
Let $X \in \Ob(\mathcal{D})$.
\begin{enumerate}
\item we say the {\it right derived functor $RF$ is defined at}
$X$ if the ind-object
$$
(X/S) \longrightarrow \mathcal{D}', \quad
(s : X \to X') \longmapsto F(X')
$$
is essentially constant\footnote{For a discussion of when an ind-object
or pro-object of a category is essentially constant we refer to
Categories, Section \ref{categories-section-essentially-constant}.};
in this case the value
$Y$ in $\mathcal{D}'$ is called the {\it value of $RF$ at $X$}.
\item we say the {\it left derived functor $LF$ is defined at} $X$
if the pro-object
$$
(S/X) \longrightarrow \mathcal{D}', \quad
(s: X' \to X) \longmapsto F(X')
$$
is essentially constant; in this case the value $Y$ in $\mathcal{D}'$
is called the {\it value of $LF$ at $X$}.
\end{enumerate}
By abuse of notation we often denote the values simply
$RF(X)$ or $LF(X)$.
\end{definition}

\noindent
It will turn out that the full subcategory of $\mathcal{D}$ consisting
of objects where $RF$ is defined is a triangulated subcategory, and
$RF$ will define a functor on this subcategory which transforms morphisms
of $S$ into isomorphisms.

\begin{lemma}
\label{lemma-derived-functor}
Assumptions and notation as in
Situation \ref{situation-derived-functor}.
Let $f : X \to Y$ be a morphism of $\mathcal{D}$.
\begin{enumerate}
\item If $RF$ is defined at $X$ and $Y$ then there exists a unique
morphism $RF(f) : RF(X) \to RF(Y)$ between the values such that
for any commutative diagram
$$
\xymatrix{
X \ar[d]_f \ar[r]_s & X' \ar[d]^{f'} \\
Y \ar[r]^{s'} & Y'
}
$$
with $s, s' \in S$ the diagram
$$
\xymatrix{
F(X) \ar[d] \ar[r] & F(X') \ar[d] \ar[r] & RF(X) \ar[d] \\
F(Y) \ar[r] & F(Y') \ar[r] & RF(Y)
}
$$
commutes.
\item If $LF$ is defined at $X$ and $Y$ then there exists a unique
morphism $LF(f) : LF(X) \to LF(Y)$ between the values such that
for any commutative diagram
$$
\xymatrix{
X' \ar[d]_{f'} \ar[r]_s & X \ar[d]^f \\
Y' \ar[r]^{s'} & Y
}
$$
with $s, s'$ in $S$ the diagram
$$
\xymatrix{
LF(X) \ar[d] \ar[r] & F(X') \ar[d] \ar[r] & F(X) \ar[d] \\
LF(Y) \ar[r] & F(Y') \ar[r] & F(Y)
}
$$
commutes.
\end{enumerate}
\end{lemma}

\begin{proof}
Part (1) holds if we only assume that the colimits
$$
RF(X) = \colim_{s : X \to X'} F(X')
\quad\text{and}\quad
RF(Y) = \colim_{s' : Y \to Y'} F(Y')
$$
exist. Namely, to give a morphism $RF(X) \to RF(Y)$ between the colimits
is the same thing as giving for each $s : X \to X'$ in $\Ob(X/S)$
a morphism $F(X') \to RF(Y)$ compatible with morphisms in the category
$X/S$. To get the morphism we choose a commutative diagram
$$
\xymatrix{
X \ar[d]_f \ar[r]_s & X' \ar[d]^{f'} \\
Y \ar[r]^{s'} & Y'
}
$$
with $s, s'$ in $S$ as is possible by MS2 and we set
$F(X') \to RF(Y)$ equal to the composition $F(X') \to F(Y') \to RF(Y)$.
To see that this is independent of the choice of the diagram above use
MS3. Details omitted. The proof of (2) is dual.
\end{proof}

\begin{lemma}
\label{lemma-derived-inverts}
Assumptions and notation as in
Situation \ref{situation-derived-functor}.
Let $s : X \to Y$ be an element of $S$.
\begin{enumerate}
\item $RF$ is defined at $X$ if and only if it is defined at $Y$.
In this case the map $RF(s) : RF(X) \to RF(Y)$ between values
is an isomorphism.
\item $LF$ is defined at $X$ if and only if it is defined at $Y$.
In this case the map $LF(s) : LF(X) \to LF(Y)$ between values
is an isomorphism.
\end{enumerate}
\end{lemma}

\begin{proof}
Omitted.
\end{proof}

\begin{lemma}
\label{lemma-derived-shift}
\begin{slogan}
Derived functors are compatible with shifts
\end{slogan}
Assumptions and notation as in
Situation \ref{situation-derived-functor}.
Let $X$ be an object of $\mathcal{D}$ and $n \in \mathbf{Z}$.
\begin{enumerate}
\item $RF$ is defined at $X$ if and only if it is defined at $X[n]$.
In this case there is a canonical isomorphism
$RF(X)[n]= RF(X[n])$ between values.
\item $LF$ is defined at $X$ if and only if it is defined at $X[n]$.
In this case there is a canonical isomorphism
$LF(X)[n] \to LF(X[n])$ between values.
\end{enumerate}
\end{lemma}

\begin{proof}
Omitted.
\end{proof}

\begin{lemma}
\label{lemma-2-out-of-3-defined}
Assumptions and notation as in
Situation \ref{situation-derived-functor}.
Let $(X, Y, Z, f, g, h)$ be a distinguished triangle of $\mathcal{D}$.
If $RF$ is defined at two out of three of $X, Y, Z$, then it is defined
at the third. Moreover, in this case
$$
(RF(X), RF(Y), RF(Z), RF(f), RF(g), RF(h))
$$
is a distinguished triangle in $\mathcal{D}'$. Similarly for $LF$.
\end{lemma}

\begin{proof}
Say $RF$ is defined at $X, Y$ with values $A, B$.
Let $RF(f) : A \to B$ be the induced morphism, see
Lemma \ref{lemma-derived-functor}.
We may choose a distinguished triangle
$(A, B, C, RF(f), b, c)$
in $\mathcal{D}'$. We claim that $C$ is a value of $RF$ at $Z$.

\medskip\noindent
To see this pick $s : X \to X'$ in $S$ such that there exists a morphism
$\alpha : A \to F(X')$ as in
Categories,
Definition \ref{categories-definition-essentially-constant-diagram}.
We may choose a commutative diagram
$$
\xymatrix{
X \ar[d]_f \ar[r]_s & X' \ar[d]^{f'} \\
Y \ar[r]^{s'} & Y'
}
$$
with $s' \in S$ by MS2. Using that $Y/S$ is filtered we can (after replacing
$s'$ by some $s'' : Y \to Y''$ in $S$) assume that there exists
a morphism $\beta : B \to F(Y')$ as in
Categories,
Definition \ref{categories-definition-essentially-constant-diagram}.
Picture
$$
\xymatrix{
A \ar[d]_{RF(f)} \ar[r]_-\alpha &
F(X') \ar[r] \ar[d]^{F(f')} &
A \ar[d]^{RF(f)} \\
B \ar[r]^-\beta & F(Y') \ar[r] & B
}
$$
It may not be true that the left square commutes, but the outer and
right squares commute.
The assumption that the ind-object $\{F(Y')\}_{s' : Y' \to Y}$
is essentially constant means that there exists a $s'' : Y \to Y''$
in $S$ and a morphism $h : Y' \to Y''$ such that $s'' = h \circ s'$ and
such that $F(h)$ equal to $F(Y') \to B \to F(Y') \to F(Y'')$. Hence
after replacing $Y'$ by $Y''$ and $\beta$ by $F(h) \circ \beta$ the
diagram will commute (by direct computation with arrows).

\medskip\noindent
Using MS6 choose a morphism of triangles
$$
(s, s', s'') : (X, Y, Z, f, g, h) \longrightarrow (X', Y', Z', f', g', h')
$$
with $s'' \in S$. By TR3 choose a morphism of triangles
$$
(\alpha, \beta, \gamma) :
(A, B, C, RF(f), b, c)
\longrightarrow
(F(X'), F(Y'), F(Z'), F(f'), F(g'), F(h'))
$$

\medskip\noindent
By Lemma \ref{lemma-derived-inverts} it suffices to prove that $RF(Z')$
is defined and that the arrow $\gamma : C \to F(Z')$ induces an isomorphism
$C \to RF(Z')$. Namely, then we will get an isomorphism
$$
(A, B, C, RF(f), b, c)
\longrightarrow
(RF(X'), RF(Y'), RF(Z'), RF(f'), RF(g'), RF(h'))
$$
of triangles and by TR1 we conclude that the target is a
distinguished triangle. Consider the category $\mathcal{I}$ of
Lemma \ref{lemma-limit-triangles}
of triangles
$$
\mathcal{I} =
\{(t, t', t'') : (X', Y', Z', f', g', h') \to (X'', Y'', Z'', f'', g'', h'')
\mid (t, t', t'') \in S\}
$$
To show that the system $F(Z'')$ is essentially constant over the category
$Z'/S$ is equivalent to showing that the system of $F(Z'')$ is essentially
constant over $\mathcal{I}$ because $\mathcal{I} \to Z'/S$ is cofinal, see
Categories, Lemma \ref{categories-lemma-cofinal-essentially-constant}
(cofinality is proven in Lemma \ref{lemma-limit-triangles}).
For any object $W$ in $\mathcal{D}'$ we
consider the diagram
$$
\xymatrix{
\colim_\mathcal{I} \Mor_{\mathcal{D}'}(W, F(X'')) \ar[d] &
\Mor_{\mathcal{D}'}(W, A) \ar[l] \ar[d] \\
\colim_\mathcal{I} \Mor_{\mathcal{D}'}(W, F(Y'')) \ar[d] &
\Mor_{\mathcal{D}'}(W, B) \ar[d] \ar[l] \\
\colim_\mathcal{I} \Mor_{\mathcal{D}'}(W, F(Z'')) \ar[d] &
\Mor_{\mathcal{D}'}(W, C) \ar[d] \ar[l] \\
\colim_\mathcal{I} \Mor_{\mathcal{D}'}(W, F(X''[1])) \ar[d] &
\Mor_{\mathcal{D}'}(W, A[1]) \ar[d] \ar[l] \\
\colim_\mathcal{I} \Mor_{\mathcal{D}'}(W, F(Y''[1])) &
\Mor_{\mathcal{D}'}(W, B[1]) \ar[l]
}
$$
where the horizontal arrows are given by composing with
$(\alpha, \beta, \gamma)$. Since filtered colimits are exact
(Algebra, Lemma \ref{algebra-lemma-directed-colimit-exact}) the left column
is an exact sequence. Thus the $5$ lemma
(Homology, Lemma \ref{homology-lemma-five-lemma}) tells us the map
$$
\colim_\mathcal{I} \Mor_{\mathcal{D}'}(W, F(Z''))
\longrightarrow
\Mor_{\mathcal{D}'}(W, C)
$$
is bijective. We conclude that $F(Z'')$ is essentially constant
over $\mathcal{I}$ with value $C$ by part (4) of
Categories, Lemma \ref{categories-lemma-characterize-essentially-constant-ind}.
\end{proof}

\begin{lemma}
\label{lemma-direct-sum-defined}
Assumptions and notation as in Situation \ref{situation-derived-functor}.
Let $X, Y$ be objects of $\mathcal{D}$.
\begin{enumerate}
\item If $RF$ is defined at $X$ and $Y$, then $RF$ is defined at $X \oplus Y$.
\item If $\mathcal{D}'$ is Karoubian and $RF$ is defined at $X \oplus Y$,
then $RF$ is defined at both $X$ and $Y$.
\end{enumerate}
In either case we have $RF(X \oplus Y) = RF(X) \oplus RF(Y)$.
Similarly for $LF$.
\end{lemma}

\begin{proof}
If $RF$ is defined at $X$ and $Y$, then the distinguished triangle
$X \to X \oplus Y \to Y \to X[1]$ (Lemma \ref{lemma-split}) and
Lemma \ref{lemma-2-out-of-3-defined}
shows that $RF$ is defined at $X \oplus Y$ and that we
have a distinguished triangle
$RF(X) \to RF(X \oplus Y) \to RF(Y) \to RF(X)[1]$.
Applying Lemma \ref{lemma-split} to this once more we find
that $RF(X \oplus Y) = RF(X) \oplus RF(Y)$.
This proves (1) and the final assertion.

\medskip\noindent
Conversely, assume that $RF$ is defined at $X \oplus Y$ and that $\mathcal{D}'$
is Karoubian. Since $S$ is a saturated system $S$ is the set of arrows which
become invertible under the additive localization functor
$Q : \mathcal{D} \to S^{-1}\mathcal{D}$, see
Categories, Lemma \ref{categories-lemma-what-gets-inverted}.
Thus for any $s : X \to X'$ and $s' : Y \to Y'$ in $S$ the morphism
$s \oplus s' : X \oplus Y \to X' \oplus Y'$ is an element of $S$.
In this way we obtain a functor
$$
X/S \times Y/S \longrightarrow (X \oplus Y)/S
$$
Recall that the categories $X/S, Y/S, (X \oplus Y)/S$ are filtered
(Categories, Remark
\ref{categories-remark-left-localization-morphisms-colimit}).
By Categories, Lemma \ref{categories-lemma-essentially-constant-over-product}
$X/S \times Y/S$ is filtered and
$F|_{X/S} : X/S \to \mathcal{D}'$ (resp.\ $G|_{Y/S} : Y/S \to \mathcal{D}'$)
is essentially constant if and only if
$F|_{X/S} \circ \text{pr}_1 : X/S \times Y/S \to \mathcal{D}'$
(resp.\ $G|_{Y/S} \circ \text{pr}_2 : X/S \times Y/S \to \mathcal{D}'$)
is essentially constant. Below we will show that the displayed functor
is cofinal, hence by
Categories, Lemma \ref{categories-lemma-cofinal-essentially-constant},
we see that $F|_{(X \oplus Y)/S}$ is essentially constant implies that
$F|_{X/S} \circ \text{pr}_1 \oplus F|_{Y/S} \circ \text{pr}_2 :
X/S \times Y/S \to \mathcal{D}'$
is essentially constant. By Homology, Lemma
\ref{homology-lemma-direct-sum-from-product-essentially-constant}
(and this is where we use that $\mathcal{D}'$ is Karoubian)
we see that 
$F|_{X/S} \circ \text{pr}_1 \oplus F|_{Y/S} \circ \text{pr}_2$
being essentially constant implies
$F|_{X/S} \circ \text{pr}_1$ and
$F|_{Y/S} \circ \text{pr}_2$ are essentially constant proving that $RF$ is
defined at $X$ and $Y$.

\medskip\noindent
Proof that the displayed functor is cofinal.
To do this pick any $t : X \oplus Y \to Z$ in $S$.
Using MS2 we can find morphisms $Z \to X'$, $Z \to Y'$
and $s : X \to X'$, $s' : Y \to Y'$ in $S$ such that
$$
\xymatrix{
X \ar[d]^s & X \oplus Y \ar[d] \ar[l] \ar[r] & Y \ar[d]_{s'} \\
X' & Z \ar[l] \ar[r] & Y'
}
$$
commutes. This proves there is a map $Z \to X' \oplus Y'$ in
$(X \oplus Y)/S$, i.e., we get part (1) of Categories, Definition
\ref{categories-definition-cofinal}. To prove part (2) it suffices
to prove that given $t : X \oplus Y \to Z$ and morphisms
$s_i \oplus s'_i : Z \to X'_i \oplus Y'_i$, $i = 1, 2$ in $(X \oplus Y)/S$
we can find morphisms $a : X'_1 \to X'$, $b : X'_2 \to X'$,
$c : Y'_1 \to Y'$, $d : Y'_2 \to Y'$ in $S$ such that
$a \circ s_1 = b \circ s_2$ and $c \circ s'_1 = d \circ s'_2$.
To do this we first choose any $X'$ and $Y'$ and maps $a, b, c, d$
in $S$; this is possible as $X/S$ and $Y/S$ are filtered. Then the
two maps $a \circ s_1, b \circ s_2 : Z \to X'$ become equal in
$S^{-1}\mathcal{D}$. Hence we can find a morphism
$X' \to X''$ in $S$ equalizing them. Similarly we find $Y' \to Y''$ in $S$
equalizing $c \circ s'_1$ and $d \circ s'_2$. Replacing $X'$ by $X''$ and
$Y'$ by $Y''$ we get $a \circ s_1 = b \circ s_2$ and
$c \circ s'_1 = d \circ s'_2$.

\medskip\noindent
The proof of the corresponding statements for $LF$ are dual.
\end{proof}

\begin{proposition}
\label{proposition-derived-functor}
Assumptions and notation as in Situation \ref{situation-derived-functor}.
\begin{enumerate}
\item The full subcategory $\mathcal{E}$ of $\mathcal{D}$ consisting of
objects at which $RF$ is defined is a strictly full triangulated
subcategory of $\mathcal{D}$.
\item We obtain an exact functor
$RF : \mathcal{E} \longrightarrow \mathcal{D}'$
of triangulated categories.
\item Elements of $S$ with either source or target
in $\mathcal{E}$ are morphisms of $\mathcal{E}$.
\item Any element of $S_\mathcal{E} = \text{Arrows}(\mathcal{E}) \cap S$
is mapped to an isomorphism by $RF$.
\item The set $S_\mathcal{E}$ is a saturated multiplicative system in
$\mathcal{E}$ compatible with the triangulated structure.
\item The functor $S_\mathcal{E}^{-1}\mathcal{E} \to S^{-1}\mathcal{D}$
is a fully faithful exact functor of triangulated categories.
\item We obtain an exact functor
$$
RF : S_\mathcal{E}^{-1}\mathcal{E} \longrightarrow \mathcal{D}'.
$$
\item If $\mathcal{D}'$ is Karoubian, then $\mathcal{E}$ is a saturated
triangulated subcategory of $\mathcal{D}$.
\end{enumerate}
A similar result holds for $LF$.
\end{proposition}

\begin{proof}
Since $S$ is saturated it contains all isomorphisms (see remark
following Categories, Definition
\ref{categories-definition-saturated-multiplicative-system}).
Hence (1) follows from Lemmas \ref{lemma-derived-inverts},
\ref{lemma-2-out-of-3-defined}, and \ref{lemma-derived-shift}.
We get (2) from Lemmas \ref{lemma-derived-functor}, \ref{lemma-derived-shift},
and \ref{lemma-2-out-of-3-defined}.
We get (3) from Lemma \ref{lemma-derived-inverts}.
Part (4) follows from Lemma \ref{lemma-derived-inverts}.
Part (5) follows from the definitions and part (3).
The fully faithfulness in (6) follows from (3) and
the definitions. The fact that
$S_\mathcal{E}^{-1}\mathcal{E} \to S^{-1}\mathcal{D}$
is exact follows from the fact that a triangle in
$S_\mathcal{E}^{-1}\mathcal{E}$ is distinguished
if and only if it is isomorphic to the image of a
distinguished triangle in $\mathcal{E}$, see proof of
Proposition \ref{proposition-construct-localization}.
The factorization of $RF : \mathcal{E} \to \mathcal{D}'$
through an exact functor $S_\mathcal{E}^{-1}\mathcal{E} \to \mathcal{D}'$
follows from Lemma \ref{lemma-universal-property-localization}.
Finally, part (8) follows from Lemma \ref{lemma-direct-sum-defined}.
\end{proof}

\noindent
Proposition \ref{proposition-derived-functor}
tells us that $RF$ lives on a maximal strictly full
triangulated subcategory of $S^{-1}\mathcal{D}$ and is an exact functor
on this triangulated category. Picture:
$$
\xymatrix{
\mathcal{D} \ar[d]_Q \ar[rrr]_F & & & \mathcal{D}' \\
S^{-1}\mathcal{D} & &
S_\mathcal{E}^{-1}\mathcal{E}
\ar[ll]_{\text{fully faithful}}^{\text{exact}} \ar[ur]_{RF}
}
$$

\begin{definition}
\label{definition-everywhere-defined}
In
Situation \ref{situation-derived-functor}.
We say $F$ is {\it right derivable}, or that $RF$ {\it everywhere defined}
if $RF$ is defined at every object of $\mathcal{D}$.
We say $F$ is {\it left derivable}, or that $LF$ {\it everywhere defined}
if $LF$ is defined at every object of $\mathcal{D}$.
\end{definition}

\noindent
In this case we obtain a right (resp.\ left) derived functor
\begin{equation}
\label{equation-everywhere}
RF : S^{-1}\mathcal{D} \longrightarrow \mathcal{D}',
\quad\text{(resp. }
LF : S^{-1}\mathcal{D} \longrightarrow \mathcal{D}'),
\end{equation}
see
Proposition \ref{proposition-derived-functor}.
In most interesting situations it is not the case that $RF \circ Q$ is
equal to $F$. In fact, it might happen that the canonical map
$F(X) \to RF(X)$ is never an isomorphism. In practice this does not happen,
because in practice we only know how to prove $F$ is right derivable by
showing that $RF$ can be computed by evaluating $F$ at judiciously chosen
objects of the triangulated category $\mathcal{D}$. This warrants
a definition.

\begin{definition}
\label{definition-computes}
In
Situation \ref{situation-derived-functor}.
\begin{enumerate}
\item An object $X$ of $\mathcal{D}$ {\it computes} $RF$ if $RF$ is defined
at $X$ and the canonical map $F(X) \to RF(X)$ is an isomorphism.
\item An object $X$ of $\mathcal{D}$ {\it computes} $LF$ if $LF$ is defined
at $X$ and the canonical map $LF(X) \to F(X)$ is an isomorphism.
\end{enumerate}
\end{definition}

\begin{lemma}
\label{lemma-computes-shift}
Assumptions and notation as in
Situation \ref{situation-derived-functor}.
Let $X$ be an object of $\mathcal{D}$ and $n \in \mathbf{Z}$.
\begin{enumerate}
\item $X$ computes $RF$ if and only if $X[n]$ computes $RF$.
\item $X$ computes $LF$ if and only if $X[n]$ computes $LF$.
\end{enumerate}
\end{lemma}

\begin{proof}
Omitted.
\end{proof}

\begin{lemma}
\label{lemma-2-out-of-3-computes}
Assumptions and notation as in
Situation \ref{situation-derived-functor}.
Let $(X, Y, Z, f, g, h)$ be a distinguished triangle of $\mathcal{D}$.
If $X, Y$ compute $RF$ then so does $Z$. Similar for $LF$.
\end{lemma}

\begin{proof}
By
Lemma \ref{lemma-2-out-of-3-defined}
we know that $RF$ is defined at $Z$ and that $RF$ applied to the
triangle produces a distinguished triangle.
Consider the morphism of distinguished triangles
$$
\xymatrix{
(F(X), F(Y), F(Z), F(f), F(g), F(h)) \ar[d] \\
(RF(X), RF(Y), RF(Z), RF(f), RF(g), RF(h))
}
$$
Two out of three maps are isomorphisms, hence so is the third.
\end{proof}

\begin{lemma}
\label{lemma-direct-sum-computes}
Assumptions and notation as in Situation \ref{situation-derived-functor}.
Let $X, Y$ be objects of $\mathcal{D}$. If $X \oplus Y$ computes $RF$, then
$X$ and $Y$ compute $RF$. Similarly for $LF$.
\end{lemma}

\begin{proof}
If $X \oplus Y$ computes $RF$, then $RF(X \oplus Y) = F(X) \oplus F(Y)$.
In the proof of Lemma \ref{lemma-direct-sum-defined} we have seen that
the functor $X/S \times Y/S \to (X \oplus Y)/S$, $(s, s') \mapsto s \oplus s'$
is cofinal. Thus by Categories, Lemma
\ref{categories-lemma-cofinal-essentially-constant} and by
characterization (4) of Categories, Lemma
\ref{categories-lemma-characterize-essentially-constant-ind}
we know that for any object $W$ in $\mathcal{D}'$ the map
$$
\Hom_{\mathcal{D}'}(F(X \oplus Y), W)
\longrightarrow
\colim_{s : X \to X', s' : Y \to Y'} \Hom_{\mathcal{D}'}(F(X' \oplus Y'), W)
$$
is bijective. Since this arrow is clearly compatible with direct sum
decompositions on both sides, we conclude that the map
$$
\Hom_{\mathcal{D}'}(F(X), W)
\longrightarrow
\colim_{s : X \to X'} \Hom_{\mathcal{D}'}(F(X'), W)
$$
is bijective (minor detail omitted). Thus by 
Categories, Lemma \ref{categories-lemma-characterize-essentially-constant-ind}
we conclude $RF$ is defined at $X$ with value $F(X)$.
Similarly for $Y$.
\end{proof}

\begin{lemma}
\label{lemma-existence-computes}
Assumptions and notation as in
Situation \ref{situation-derived-functor}.
\begin{enumerate}
\item If for every object $X \in \Ob(\mathcal{D})$
there exists an arrow $s : X \to X'$ in $S$ such that $X'$ computes
$RF$, then $RF$ is everywhere defined.
\item If for every object $X \in \Ob(\mathcal{D})$
there exists an arrow $s : X' \to X$ in $S$ such that $X'$ computes
$LF$, then $LF$ is everywhere defined.
\end{enumerate}
\end{lemma}

\begin{proof}
This is clear from the definitions.
\end{proof}

\begin{lemma}
\label{lemma-find-existence-computes}
Assumptions and notation as in
Situation \ref{situation-derived-functor}.
If there exists a subset $\mathcal{I} \subset \Ob(\mathcal{D})$
such that
\begin{enumerate}
\item for all $X \in \Ob(\mathcal{D})$
there exists $s : X \to X'$ in $S$ with $X' \in \mathcal{I}$,
and
\item for every arrow $s : X \to X'$ in $S$ with $X, X' \in \mathcal{I}$
the map $F(s) : F(X) \to F(X')$ is an isomorphism,
\end{enumerate}
then $RF$ is everywhere defined and every $X \in \mathcal{I}$
computes $RF$. Dually, if there exists a subset
$\mathcal{P} \subset \Ob(\mathcal{D})$
such that
\begin{enumerate}
\item for all $X \in \Ob(\mathcal{D})$
there exists $s : X' \to X$ in $S$ with $X' \in \mathcal{P}$,
and
\item for every arrow $s : X \to X'$ in $S$ with $X, X' \in \mathcal{P}$
the map $F(s) : F(X) \to F(X')$ is an isomorphism,
\end{enumerate}
then $LF$ is everywhere defined and every $X \in \mathcal{P}$
computes $LF$.
\end{lemma}

\begin{proof}
Let $X$ be an object of $\mathcal{D}$.
Assumption (1) implies that the arrows $s : X \to X'$ in $S$ with
$X' \in \mathcal{I}$ are cofinal in the category $X/S$. Assumption
(2) implies that $F$ is constant on this cofinal subcategory.
Clearly this implies that $F : (X/S) \to \mathcal{D}'$ is essentially
constant with value $F(X')$ for any $s : X \to X'$ in $S$
with $X' \in \mathcal{I}$.
\end{proof}

\begin{lemma}
\label{lemma-compose-derived-functors-general}
Let $\mathcal{A}, \mathcal{B}, \mathcal{C}$ be triangulated categories.
Let $S$, resp.\ $S'$ be a saturated multiplicative system in
$\mathcal{A}$, resp.\ $\mathcal{B}$ compatible with the triangulated structure.
Let $F : \mathcal{A} \to \mathcal{B}$ and $G : \mathcal{B} \to \mathcal{C}$
be exact functors. Denote $F' : \mathcal{A} \to (S')^{-1}\mathcal{B}$ the
composition of $F$ with the localization functor.
\begin{enumerate}
\item If $RF'$, $RG$, $R(G \circ F)$ are everywhere defined, then there
is a canonical transformation of functors
$t : R(G \circ F) \longrightarrow RG \circ RF'$.
\item If $LF'$, $LG$, $L(G \circ F)$ are everywhere defined, then there
is a canonical transformation of functors
$t : LG \circ LF' \to L(G \circ F)$.
\end{enumerate}
\end{lemma}

\begin{proof}
In this proof we try to be careful. Hence let us think of
the derived functors as the functors
$$
RF' : S^{-1}\mathcal{A} \to (S')^{-1}\mathcal{B}, \quad
R(G \circ F) : S^{-1}\mathcal{A} \to \mathcal{C}, \quad
RG : (S')^{-1}\mathcal{B} \to \mathcal{C}.
$$
Let us denote
$Q_A : \mathcal{A} \to S^{-1}\mathcal{A}$ and
$Q_B : \mathcal{B} \to (S')^{-1}\mathcal{B}$
the localization functors. Then $F' = Q_B \circ F$. Note that for
every object $Y$ of $\mathcal{B}$ there is a canonical map
$$
G(Y) \longrightarrow RG(Q_B(Y))
$$
in other words, there is a transformation of functors
$t' : G \to RG \circ Q_B$. Let $X$ be an object of $\mathcal{A}$.
We have
\begin{align*}
R(G \circ F)(Q_A(X))
& = \colim_{s : X \to X' \in S} G(F(X')) \\
& \xrightarrow{t'} \colim_{s : X \to X' \in S} RG(Q_B(F(X'))) \\
& = \colim_{s : X \to X' \in S} RG(F'(X')) \\
& = RG(\colim_{s : X \to X' \in S} F'(X')) \\
& = RG(RF'(X)).
\end{align*}
The system $F'(X')$ is essentially constant in the category
$(S')^{-1}\mathcal{B}$. Hence we may pull the colimit inside the
functor $RG$ in the third equality of the diagram above, see
Categories, Lemma \ref{categories-lemma-image-essentially-constant}
and its proof. We omit the proof this defines a transformation
of functors. The case of left derived functors is similar.
\end{proof}




